% Source-aligned English for printed p.144, with logged editorial qualification.
\SourcePage{144}

\LayerTarget{tohoku:construction:derived-bifunctor:p2}
\noindent right partial derived functors of $T$ with respect to the second
variable $B$,
\[
 R_2^iT(A,B)=H^i(T(A,C(B)))
\]
where $C(B)$ is the complex defined by a right resolution of $B$ by injective
objects. Of course, the $R_2^iT$ are bifunctors.
\LayerTarget{tohoku:construction:derived-bifunctor:cohomological-first-variable}
\textit{Suppose now that for every injective object $B$ of $\mathbf{C}_2$, the functor
$A\to T(A,B)$ on $\mathbf{C}_1$ is exact.} We shall show that \textit{for every}
$B\in \mathbf{C}_2$ the system $(R_2^iT(A,B))$ can be regarded as a cohomological
functor of $A$. Let $C(B)$ be the complex defined by a right resolution of
$B$ by injective objects. For every exact sequence
$0\to A'\to A\to A''\to0$
in $\mathbf{C}_1$, the sequence of complex homomorphisms
$0\to T(A'',C(B))\to T(A,C(B))
 \to T(A',C(B))\to0$
is exact, since the terms of $C(B)$ are injective and by the hypothesis on
$T$. It therefore defines a long exact cohomology sequence
\LayerTarget{tohoku:sequence:partial-derived-long-exact}
\[
\ldotp\,\ldotp\,\ldotp\,\ldotp\to R_2^iT(A'',B)\to R_2^iT(A,B)
\to R_2^iT(A',B)
\mathrel{\mathop{\longrightarrow}\limits^{\partial}}R_2^{i+1}T(A'',B)
\to\ldotp\,\ldotp\,\ldotp\,\ldotp .
\]
One immediately checks that the morphism $\partial$ in this sequence is independent of the particular
complex $C(B)$ chosen and commutes with homomorphisms of exact sequences.
Thus $(R_2^iT(A,B))$, for fixed $B$, is indeed a cohomological functor of $A$.
Moreover, one immediately checks that if $B\to B'$ is a morphism in $\mathbf{C}_2$, the corresponding
morphisms $R_2^iT(A,B)\to R_2^iT(A,B')$ define a morphism of
$\partial$-functors in $A$.
\LayerTarget{tohoku:def:right-balanced}
Suppose also that every object of $\mathbf{C}_1$ is isomorphic to a quotient of a
projective object and that, for every projective $A\in \mathbf{C}_1$, the functor
$T(A,B)$ is exact in $B$. In the terminology of [6], $T$ is then
``right balanced.'' The $R_2^iT(A,B)$ are then also the right derived functors
$R^iT(A,B)$ of $T$, and they also coincide with the partial derived functors
$R_1^iT(A,B)$. In this case the boundary homomorphisms constructed above are
precisely the natural boundary homomorphisms of $R^iT$ and $R_1^iT$ with
respect to their first variable.

\LayerTarget{tohoku:example:ext-bifunctor}
The preceding discussion was developed chiefly for an abelian category $\mathbf{C}$
in which every object is isomorphic to a subobject of an injective object,
without assuming that every object is isomorphic to a quotient of a projective
object. Take $\mathbf{C}_1=\mathbf{C}_2=\mathbf{C}$, let $C'$ be the category of abelian groups, and set
$T(A,B)=\mathrm{Hom}(A,B)$. By the definition of an injective object,
$\mathrm{Hom}(A,B)$ is exact in $A$ when $B$ is injective. Denote the right
derived functors with respect to $B$ by $\mathrm{Ext}^p(A,B)$. The system
$(\mathrm{Ext}^p(A,B))$ is a cohomological functor both in $B$ \textit{and} in $A$.
Since $\mathrm{Hom}(A,B)$ is left exact,
$\mathrm{Ext}^0(A,B)=\mathrm{Hom}(A,B)$; it is also easy to see that
$\mathrm{Ext}^1(A,B)$ can be interpreted as the group of classes of extensions
of $A$ as quotient by $B$ as subobject; see [6, Chap. XIV].

\LayerTarget{tohoku:example:sheaf-ext}
An important class of cases in which an abelian category $\mathbf{C}$ has enough
injective objects and may fail to have enough projective objects is furnished
by categories of sheaves of modules over a given sheaf of rings on a
topological space $X$. In Chapter IV we shall study the groups
$\mathrm{Ext}^p(A,B)$ in this case in greater detail.

\LayerTarget{tohoku:section:2.4}
\textbf{2.4. Spectral sequences and spectral functors.} For the
theory of

\TohokuContinuation{145}
